%%%%%%%%%%%%%%%%%
% General CV - Guillaume Boileau
%%%%%%%%%%%%%%%%%

\documentclass[8pt,a4paper,ragged2e,withhyper]{../_shared/altacv}

\geometry{left=1.25cm,right=1.25cm,top=.5cm,bottom=1.cm,columnsep=1.2cm}

\usepackage{paracol}
\usepackage{fontawesome5}
\usepackage{hyperref}

\iftutex 
  \setmainfont{Roboto Slab}
  \setsansfont{Lato}
  \renewcommand{\familydefault}{\sfdefault}
\else
  \usepackage[rm]{roboto}
  \usepackage[defaultsans]{lato}
  \renewcommand{\familydefault}{\sfdefault}
\fi

\definecolor{SlateGrey}{HTML}{2E2E2E}
\definecolor{LightGrey}{HTML}{666666}
\definecolor{DarkPastelRed}{HTML}{8AC0D9}
\definecolor{PastelRed}{HTML}{00B0DA}
\definecolor{GoldenEarth}{HTML}{B4AD0C}

\colorlet{name}{black}
\colorlet{tagline}{PastelRed}
\colorlet{heading}{DarkPastelRed}
\colorlet{headingrule}{GoldenEarth}
\colorlet{subheading}{PastelRed}
\colorlet{accent}{PastelRed}
\colorlet{emphasis}{SlateGrey}
\colorlet{body}{LightGrey}

\definecolor{violet}{HTML}{5A2582}
\definecolor{rouge}{HTML}{F53B3B}
\definecolor{noir}{HTML}{00232C}
\definecolor{bleu}{HTML}{00B0DA}
\definecolor{rose}{HTML}{F831A0}
\definecolor{jaune}{HTML}{FFEC00}
\definecolor{vert}{HTML}{34AD6C}

\renewcommand{\namefont}{\Huge\rmfamily\bfseries}
\renewcommand{\personalinfofont}{\footnotesize}
\renewcommand{\cvsectionfont}{\LARGE\rmfamily\bfseries}
\renewcommand{\cvsubsectionfont}{\large\bfseries}
\renewcommand{\cvItemMarker}{{\small\textbullet}}
\renewcommand{\cvRatingMarker}{\faCircle}

\input{../_shared/pubs-num.tex}
\addbibresource{../_shared/sample.bib}

\begin{document}

\name{BOILEAU Guillaume}
\tagline{Ingénieur docteur | Data, logiciel, systèmes complexes, validation \& coordination technique}

\photoL{2.8cm}{../_shared/moi}

\personalinfo{%
  \email{guillaume.boileaupro@gmail.com}
  \phone{+33 6 59 94 85 84}
  \location{Alpes-Maritimes, FRANCE}
  \linkedin{guillaume-boileau-891b81265}
  \researchgate{Guillaume-Boileau-2}
  \orcid{0000-0002-3576-6968}
}

\makecvheader
\vspace{-0.3cm}

\AtBeginEnvironment{itemize}{\small}
\columnratio{0.64}

\begin{paracol}{2}

\cvsection{Experience}

\cvevent{\textbf{Ingénieur logiciel full stack}}{VEV Platform Services France}{2025 -- 2026}{Valbonne, France}

\textbf{Contexte :} Développement logiciel et support opérationnel pour une plateforme CPMS liée aux infrastructures de recharge de véhicules électriques.

\textbf{Objectif :} Contribuer au développement, à l'intégration et au suivi technique de services logiciels connectés à des équipements terrain et à des flux de données opérationnels.

\begin{itemize}
  \item \textbf{Développement logiciel :} Contribution au développement d'applications web et de services backend en TypeScript, Angular et Node.js.
  \item \textbf{Intégration système :} Travail sur des interfaces entre services logiciels, données opérationnelles et infrastructures de recharge.
  \item \textbf{Flux de données :} Suivi et analyse de flux opérationnels, notamment via des échanges FTP utilisés pour le pilotage et le suivi technique.
  \item \textbf{Support production :} Analyse d'incidents, investigation de comportements applicatifs et contribution à la fiabilisation des services.
  \item \textbf{Qualité et cohérence :} Vérification de la cohérence entre données, services applicatifs, comportement terrain et attentes opérationnelles.
  \item \textbf{Organisation agile :} Utilisation de Zenhub pour le suivi des tâches, la priorisation, le backlog et la coordination dans un environnement de développement itératif.
  \item \textbf{Communication technique :} Formalisation d'observations, d'analyses d'incidents et de points techniques pour faciliter la prise de décision.
\end{itemize}

\divider

\cvevent{\textbf{Ingénieur R\&D senior -- Data, logiciel \& validation}}{Laboratoire Lagrange, Observatoire de la Côte d'Azur, CNES}{2023 -- 2025}{Nice, France}

\textbf{Contexte :} Activités d'ingénierie et de R\&D pour la mission spatiale LISA ESA/NASA, avec des enjeux de simulation, de traitement de données, d'intégration de modèles et de validation technique.

\textbf{Objectif :} Développer des outils Python permettant de structurer, comparer, intégrer et valider des modèles et des résultats issus de contributions hétérogènes.

\begin{itemize}
  \item \textbf{Développement Python :} Conception et développement de workflows pour l'intégration, la comparaison et la validation de modèles.
  \item \textbf{Pipelines de données :} Mise en place de chaînes de traitement pour organiser les entrées, transformer les sorties et comparer les résultats.
  \item \textbf{Validation technique :} Définition de contrôles de cohérence, de règles de comparaison et de méthodes de vérification des résultats.
  \item \textbf{Intégration de modèles :} Consolidation de résultats produits par différents modèles dans un cadre commun d'analyse et de revue.
  \item \textbf{Traçabilité :} Structuration des configurations, sorties de simulation, résultats d'analyse et éléments de documentation.
  \item \textbf{Documentation et capitalisation :} Utilisation de Confluence pour organiser la documentation technique, le suivi de validation, la traçabilité des analyses et le partage de connaissances.
  \item \textbf{Coordination :} Travail avec des contributeurs scientifiques, logiciels et institutionnels dans un environnement international.
  \item \textbf{Plateforme logicielle :} Développement de CATpyLISA, plateforme Python pour l'intégration, la comparaison, la validation et la revue technique de modèles.
  \textcolor{DarkPastelRed}{\href{https://gitlab.oca.eu/gboileau/lisa_l3_tools}{\bf GitLab:CATpyLISA}}
\end{itemize}

\divider

\cvevent{\textbf{Data Scientist -- Modélisation \& analyse de données}}{Universiteit Antwerpen, Belgium}{2021 -- 2023}{Antwerp, Belgium}

\textbf{Contexte :} Analyse de données et études de performance pour des instruments de précision dans un environnement technique international.

\textbf{Objectif :} Développer des workflows robustes et reproductibles pour identifier, caractériser et réduire des sources affectant la qualité de mesure.

\begin{itemize}
  \item \textbf{Analyse de données :} Traitement, structuration et analyse de données complexes issues d'instruments de précision.
  \item \textbf{Modélisation statistique :} Développement de méthodes d'analyse et d'inférence bayésienne, notamment par chaînes MCMC.
  \item \textbf{Traitement du signal :} Identification et caractérisation de contributions de bruit dans des mesures complexes.
  \item \textbf{Pipelines Python :} Mise en place de workflows de traitement, filtrage, transformation et analyse statistique.
  \item \textbf{Validation de résultats :} Analyse de sensibilité, vérification de cohérence et évaluation des limites de performance.
  \item \textbf{Reporting technique :} Production de résultats traçables, de synthèses techniques et de supports pour des parties prenantes internationales.
\end{itemize}

\divider

\cvevent{\textbf{Ingénieur R\&D junior -- Signal, simulation \& analyse de données}}{ARTEMIS, Observatoire de la Côte d'Azur}{2018 -- 2021}{Nice, France}

\textbf{Contexte :} Activités de R\&D liées à l'analyse de signaux faibles, à la simulation numérique et à la préparation de missions spatiales et d'instruments de détection.

\textbf{Objectif :} Développer des méthodes d'analyse, automatiser des simulations et produire des résultats exploitables dans un environnement scientifique et technique exigeant.

\begin{itemize}
  \item \textbf{Simulation numérique :} Développement d'approches de simulation pour des environnements de signaux complexes et des jeux de données de grande taille.
  \item \textbf{Analyse de données :} Mise en place de méthodes pour analyser et séparer des signaux faibles et superposés.
  \item \textbf{Traitement du signal :} Études de caractérisation, de décomposition et de comparaison de signaux dans des contextes bruités.
  \item \textbf{Automatisation :} Développement de workflows reproductibles pour les analyses répétées, les comparaisons et les études de sensibilité.
  \item \textbf{Investigation technique :} Analyse de sources de bruit et réalisation de diagnostics par corrélations entre capteurs.
  \item \textbf{Documentation :} Formalisation des méthodes, hypothèses, résultats et conclusions pour revue collaborative.
\end{itemize}

\switchcolumn

\cvsection{Profile}

\textbf{Ingénieur docteur avec une double compétence en data, logiciel et systèmes complexes, spécialisé dans l'analyse de données, le développement de pipelines, la validation technique et la coordination de projets multidisciplinaires.}

Mon parcours combine développement Python, traitement de données, simulation numérique, intégration de modèles, investigation d'anomalies, documentation technique et support à la décision. J'ai travaillé dans des environnements exigeants : mission spatiale LISA ESA/NASA, instrumentation scientifique, plateformes logicielles industrielles et infrastructures de recharge électrique.

\medskip

\textbf{Positionnement cible :} ingénieur systèmes / IVVQ, chef de projet technique, PMO technique, data / AI engineer, backend Python, coordinateur R\&T ou documentation technique.

\begin{itemize}
  \item Capacité à faire le lien entre besoins métier, contraintes techniques, données, logiciel, qualité et livraison.
  \item Expérience en développement logiciel, data processing, simulation, validation de résultats et automatisation de workflows.
  \item Culture forte de la qualité : cohérence, traçabilité, reproductibilité, documentation et contrôle des résultats.
  \item Habitude des environnements collaboratifs, internationaux, multidisciplinaires et orientés livraison.
\end{itemize}

\cvsection{Languages}

\cvskill{English}{4}
\divider
\cvskill{Spanish}{2}
\divider

\medskip

\cvsection{Education}

\cvevent{PhD in Physics}{Université Côte d'Azur}{2018 -- 2021}{Nice, France}

\divider

\cvevent{Selective Graduate Program in Fundamental Physics (Magistère)}{Université Paris-Saclay}{2016 -- 2018}{Orsay, France}

\divider

\cvevent{MSc in Astronomy and Astrophysics}{Observatoire de Paris / Université Paris-Saclay}{2017 -- 2018}{Paris, France}

\divider

\cvevent{BSc in Fundamental Physics}{Université Paris-Saclay}{2015 -- 2016}{Orsay, France}

\cvsection{Professional Training}

\cvevent{Supervised Machine Learning: Regression \& Classification}{DeepLearning.AI -- Andrew Ng}{2026}{Coursera}
\divider

\cvevent{Développement assisté par IA}{Prompting, relecture de code, debugging, prompting incrémental et critique}{2026}{Autoformation}

\end{paracol}

\cvsection{Projects}

\cvevent{CATpyLISA -- Plateforme Python d'intégration, comparaison et validation}{Mission spatiale LISA ESA/NASA}{2023 -- 2025}{}

Développement d'une plateforme Python dédiée à l'intégration de modèles, à la structuration de données, à la comparaison de résultats, à la validation technique et à la revue de workflows liés à la mission LISA.

\textbf{Rôle :} développement Python, conception de pipelines, structuration de données, logique de validation, contrôles de cohérence, documentation technique et coordination avec plusieurs contributeurs.

\textbf{Compétences mobilisées :} Python, data processing, automatisation, modélisation, validation, documentation, GitLab, environnement collaboratif international.

\divider

\cvevent{Flux de données opérationnels pour infrastructures de recharge}{Plateforme CPMS -- VEV Platform Services France}{2025 -- 2026}{}

Contribution à des services logiciels liés aux échanges de données opérationnelles et au comportement d'infrastructures connectées de recharge de véhicules électriques.

\textbf{Rôle :} développement logiciel, suivi de flux de données, analyse d'incidents, diagnostic opérationnel, fiabilisation d'interfaces et coordination technique.

\textbf{Compétences mobilisées :} TypeScript, Angular, Node.js, FTP, API, intégration logiciel/système, support production, analyse d'anomalies, Zenhub.

\divider

\cvevent{Collaborations scientifiques et techniques internationales}{LISA, LIGO--Virgo--KAGRA, Einstein Telescope}{2018 -- now}{}

Contribution à des travaux collaboratifs en analyse de données, traitement du signal, simulation, caractérisation de bruit, validation de résultats et investigation technique dans des environnements internationaux.

\textbf{Rôle :} développement d'outils d'analyse, contribution à des workflows de traitement, études techniques, documentation et revue collaborative.

\textbf{Compétences mobilisées :} analyse de données, traitement du signal, modélisation, simulation, validation, coordination technique et communication scientifique.

\cvsection{Compétences clés}

\vspace{-.3cm}

\cvsubsection{Pilotage technique, projet \& coordination}

{\small
\cvtag{Coordination technique}
\cvtag{Gestion de projet technique}
\cvtag{PMO technique}
\cvtag{Suivi de tâches}
\cvtag{Priorisation}
\cvtag{Reporting}
\cvtag{Gestion des risques}
\cvtag{Stakeholder management}
\cvtag{Agile}
\cvtag{Kanban}
\cvtag{Cycle en V}
\cvtag{Confluence}
\cvtag{Jira}
\cvtag{Zenhub}
}

\divider\smallskip
\vspace{-.3cm}

\cvsubsection{IVVQ, validation, qualité \& systèmes}

{\small
\cvtag{IVVQ logiciel}
\cvtag{Vérification / Validation}
\cvtag{Intégration système}
\cvtag{Tests fonctionnels}
\cvtag{Non-régression}
\cvtag{Traçabilité exigences}
\cvtag{Contrôles de cohérence}
\newline
\cvtag{Analyse d'anomalies}
\cvtag{Diagnostic technique}
\cvtag{Root cause analysis}
\cvtag{Documentation qualité}
\cvtag{Analyse de risques}
\cvtag{Systèmes complexes}
}

\divider\smallskip
\vspace{-.3cm}

\cvsubsection{Data, IA, modélisation \& signal}

{\small
\cvtag{Analyse de données}
\cvtag{Data processing}
\cvtag{Data quality}
\cvtag{Pipelines de données}
\cvtag{Modélisation statistique}
\cvtag{Machine Learning}
\cvtag{Classification}
\cvtag{Scikit-learn}
\cvtag{PyTorch}
\cvtag{MCMC}
\cvtag{Inférence bayésienne}
\cvtag{Simulation numérique}
\cvtag{Traitement du signal}
\cvtag{Analyse de bruit}
\cvtag{Signaux faibles}
\cvtag{Prompt engineering}
\cvtag{LLM awareness}
}

\divider\smallskip
\vspace{-.3cm}

\cvsubsection{Développement logiciel, web \& automatisation}

{\small
\cvtag{Python}
\cvtag{Pandas}
\cvtag{NumPy}
\cvtag{SciPy}
\cvtag{SQL}
\cvtag{TypeScript}
\cvtag{Angular}
\cvtag{Node.js}
\cvtag{REST API}
\cvtag{Bash}
\cvtag{Linux}
\cvtag{Git}
\cvtag{GitLab}
\cvtag{Docker}
\cvtag{CI/CD}
\cvtag{FTP}
\cvtag{Power BI}
\cvtag{n8n awareness}
\cvtag{HubSpot awareness}
}

\divider\smallskip
\vspace{-.3cm}

\cvsubsection{Domaines d'application}

{\small
\cvtag{Spatial}
\cvtag{LISA ESA/NASA}
\cvtag{Systèmes satellites}
\cvtag{Instrumentation scientifique}
\cvtag{Plateformes data}
\cvtag{CPMS}
\cvtag{Recharge électrique}
\cvtag{Systèmes industriels}
\newline
\cvtag{Projets internationaux}
}

\end{document}